\chapter{Introduksjon}

Helsetjenestens driftsorganisasjon for nødnett HF (HDO) leverer kommunikasjonsløsninger til den akuttmedisinske kjeden i Norge. En del av dette arbeidet er videoløsninger som brukes av akuttmedisinske kommunikasjonssentraler (AMK) og legevaktsentraler (LVS). Når en innringer kontakter en sentral, kan operatøren sende en SMS-lenke som lar innringeren dele video fra mobilkameraet sitt. Lyden går fortsatt gjennom telefonsamtalen, mens videoen brukes som støtte for å forstå situasjonen bedre.

HDO bruker i dag to ulike videomotorer: AntMedia Server og NHN Video. Begge kan brukes til å opprette og håndtere videostrømmer, men de har ulike API-er, autentiseringsløsninger og måter å representere videorom på. Det gjør at ny funksjonalitet må tilpasses hver motor separat. Over tid kan dette gi mer vedlikehold, mer dobbeltarbeid og sterkere kobling til bestemte leverandører.

Målet med denne bacheloroppgaven er å utvikle og vurdere en mellomvare som skjuler disse forskjellene bak ett felles API. I stedet for at klientsystemene må forholde seg direkte til AntMedia eller NHN Video, skal de kunne sende forespørsler til VideoAPI. Mellomvaren velger riktig videomotor, oversetter forespørselen til riktig format og returnerer et felles svar tilbake til klienten.

Oppgaven undersøker derfor følgende forskningsspørsmål:

\begin{quote}
I hvilken grad kan en provider-basert mellomvarearkitektur standardisere tilgang til AntMedia Server og NHN Video gjennom ett felles API, samtidig som backend-spesifikk logikk isoleres fra kjernelogikken i VideoAPI?
\end{quote}

For å svare på dette ble det utviklet en proof-of-concept-løsning i ASP.NET Core. Løsningen inneholder et REST-API for håndtering av videorom, egne providere for AntMedia Server og NHN Video, tokenbasert autentisering, Docker-støtte, helsesjekker, strukturert logging og Prometheus-metrikker. Arbeidet ble evaluert gjennom kravsporing, automatiserte tester, ende-til-ende-testing mot HDOs testmiljøer og tilbakemelding fra HDO.

Rapporten er strukturert slik: Kapittel 2 gir bakgrunn om HDO, videoløsningene og relevante tekniske konsepter. Kapittel 3 beskriver kravene fra HDO og hvordan de er tolket i prosjektet. Kapittel 4 beskriver arbeidsprosessen og metoden som ble brukt. Kapittel 5 forklarer arkitekturen og de viktigste designvalgene. Kapittel 6 beskriver implementasjonen. Kapittel 7 presenterer testing og evaluering. Kapittel 8 diskuterer resultatene, begrensningene, bærekraft, bruk av kunstig intelligens og videre arbeid.